\section{中国法制史}
\begin{enumerate}
  \item 不要占太多时间。分值比例非常低。只关注中法史。
  \item 全选一个选项，蒙 1 分。
  \item 西周。以\zd{德}配天，明德慎罚。制\zd{礼}作乐。
        \begin{enumerate}
          \item 礼就是法。一体两面。
          \item \zd{出礼则入刑} vs 礼刑法。尊尊君为首，亲亲父为首。
          \item 刑不上大夫，礼不下庶人。强调等级，不是不受刑，\qd{死法不一样}。
          \item 五刑。墨、劓（鼻）、剕（脚）、宫、大辟（各种死刑的统称）。
          \item 秘密法。法律不公开。
          \item 眚（瞎）\yjt 过失，非眚 \yjt 故意；惟终 \yjt 自始自终，非终 \yjt 并非自始自终。
          \item \zd{质}活\zd{剂}死。买卖合同。制人。
          \item 傅别。傅书于上，别则为二。借贷。
          \item 婚姻家庭继承。一直到清朝，没有本质变化。
                \begin{enumerate}
                  \item 一夫一妻；同姓不婚；父母之命媒妁之言。
                  \item 六礼。纳采（提亲）、问名（生辰八字）、纳吉、纳征（彩礼）、请期（婚期）、亲迎。
                  \item 解除（男方单向）。七出三不去。前贫后富、守孝 3 年、娘家没人。
                  \item 嫡长子继承。立子以贵（妻）不以长（妾），立嫡以长不以贤。
                \end{enumerate}
          \item 司寇。
          \item 狱（刑事）讼（民事）。
          \item 五听三刺。辞、色、气、耳、目；收集、听取意见，群臣 \yjt 官吏 \yjt 国人。\qd{能面刺寡人者}。
          \item 三赦。幼弱、老耄、蠢愚者。
        \end{enumerate}
  \item 春秋战国。
        \begin{enumerate}
          \item \zd{公布成文法}。铸刑书，铸刑鼎。书前鼎后。打破刑不可知，威不可测。
          \item 李悝《法经》。轻罪重刑，明法重刑。
                \begin{enumerate}
                  \item 第一部比较系统的成文法典。
                  \item 六篇。先分后总。程序实体不分。
                        \begin{enumerate}
                          \item 盗。财产。
                          \item 贼安全。
                          \item 网。囚禁、审判。
                          \item 捕。
                          \item 杂六禁。淫、狡、城、嬉（赌博）、徒（集会、游行）。
                          \item 具（总则）。加减。
                        \end{enumerate}
                  \item 源头 \yjt
                        商鞅 \yjt
                        汉萧何，增加户（户口）、兴（城市建设）厩（养牛羊马），九章律 \yjt
                        魏晋南北朝，新律 \yjt
                        两晋 \yjt
                        北朝，北齐律 \yjt
                        唐宋明清
                \end{enumerate}
          \item 商鞅变法。改法为律。强调统一适用，普遍调整；通过法律富国强兵，奖励军功。
        \end{enumerate}
  \item 秦。刑罚适用原则、罪名部分。1 道小题目。\zd{以法为教，以吏为师}。
        \begin{enumerate}
          \item 危害皇权类。谋反、偶语诗书、以古非今、诅咒、诽谤、妄言、妖言、非所宜言、投书（匿名信）、不行军令。
          \item 侵犯财产和人身。共盗（5 人）、群盗（聚众反抗统治秩序）。
          \item 贼杀。特别严重德杀人。
          \item 渎职。见知（犯罪）不举。不直（故意犯罪）、失刑、纵囚。
          \item 妨害管理。违令卖酒、逋事（逃避徭役）、乏徭、
          \item 婚姻家庭。夫妻、夫通、奸妻、私逃、夫殴妻，夫通奸。擅杀子、子不孝、子女控告父母。卑幼殴尊长、乱伦。
          \item 笞刑（主刑）；徒刑（副刑）。修城、旦、舂米、鬼薪白粲、隶臣妾。
          \item 按身高确定刑事能力。
        \end{enumerate}
  \item 汉。儒家化。\zd{汉承秦制，德主刑辅，礼刑并用}。
        \begin{enumerate}
          \item 文景废肉刑。
          \item 上请（减轻刑罚）。恤刑。亲亲得相首匿。
          \item 《春秋》决狱，秋冬行刑（死刑）。核心看动机。
                \begin{enumerate}
                  \item 本其事而原其志。
                  \item 邪恶者不待成 \yjt 未遂既遂。
                  \item 首恶者罪特重 \yjt 首犯从犯。
                  \item 本直者其论轻 \yjt 看重主观动机。
                \end{enumerate}
          \item 廷尉。御史。
        \end{enumerate}
  \item 魏晋南北朝。注重具体的技术操作。\zd{律政分途，律重训诂}。
        \begin{enumerate}
          \item 法律解释。张斐、杜预。故杀、斗杀、过失杀 \yjt 唐律七杀理论。
          \item 法律形式和法典结构、内容的变化。
                \begin{enumerate}
                  \item 多种形式相互为用，律令科（补充变通）比（判例）格（补充律，刑事性质）式交相为用。
                  \item 先总后分。
                \end{enumerate}
          \item 曹魏律。
                \begin{enumerate}
                  \item 具法，刑名置于律首。
                  \item 八议入律。亲故贤能，功贵勤宾，可以商量着办。
                \end{enumerate}
          \item 晋律、北齐律。准五服以制罪。解决亲属犯罪定罪问题。
          \item 北魏律。官当入律，死刑复核。增加鞭刑、杖刑。
          \item 北齐律。名例律（总则）。重罪 10 条。
          \item 开皇律。确立了封建五刑，笞、杖、徒、流、死。
          \item 大理寺（审判）、御史台（监察）。
        \end{enumerate}
  \item 隋。开皇律。立法宽松简约；十罪改十恶。
  \item 唐律和中华法系。必有 1 题。\zd{礼法合一、礼律合一、礼刑合一}。相对完整的复习。
        \begin{enumerate}
          \item 唐律疏议。\zd{中华法系}的立法标杆。
          \item 唐六典。户、礼、吏、兵、刑、工。行政法典。规定了\zd{回避}制度。
          \item 六赃，六杀，保辜 \yjt 伤情观察期。
          \item 区分公罪、私罪。公职人员。
          \item 自首。
          \item 类推原则。举重以明轻。
          \item 化外。涉外。
          \item 大理寺、刑部（复核）、御史台。
          \item 三司推事。各派一个。
          \item 都堂集议。
          \item 先以情，再以礼，反复参验。
          \item 据状断之。众证定罪。
        \end{enumerate}
  \item 宋元。第一部刊印发行的法典。宋刑统。
        \begin{enumerate}
          \item 编敕（皇帝发布的）。
          \item 法深无善治。
          \item 折杖刑。
          \item 民法部分完整看一下。典卖。
                \begin{enumerate}
                  \item 离婚。夫外出 3 年不归，6 年不通问。夫亡而妻不守。恩断义绝。
                  \item 诸子均分。
                \end{enumerate}
          \item 翻异别堪。翻供换人。
        \end{enumerate}
  \item 明。提高了刑的地位。重典治世。\zd{明刑弼教}。
        \begin{enumerate}
          \item 奸党罪。
        \end{enumerate}
  \item 清朝。大清律例。
        \begin{enumerate}
          \item 事例。皇帝发布的 \yjt（整理）成例。则例（行政），条例（刑法）。单行法。
          \item 会审制度。
          \item 国体之争。
          \item 预备立宪。
          \item 清末律政。
          \item 领事裁判权。治外法权。
        \end{enumerate}
  \item 中华民国。
        \begin{enumerate}
          \item 三民主义。联俄联共扶助农工新旧分水岭。民族、民权、民生。
          \item 五权宪法。增加考试、监察。
          \item 五院制。
          \item 章太炎、宋教仁。
        \end{enumerate}
  \item 南京临时政府。临时政府大纲。临时约法。建立起了律师制度。
  \item 北京政府。中华民国约法。复辟。袁世凯、曹锟。
  \item 南京国民政府。蒋介石。六法全书。训政纲领。新刑罚。一告九不理。
  \item 中国共产党民主政权宪法性文件。
        \begin{enumerate}
          \item 苏维埃宪法大纲。
          \item 陕甘宁边区施政纲领。普及文化教育、普选制、三三制、男女同权。
          \item 陕甘宁边区宪法原则。
        \end{enumerate}
\end{enumerate}